\documentclass[11pt]{article}
\usepackage[margin=1in]{geometry}
\usepackage{amsmath,amssymb,amsthm,hyperref}
\newtheorem{proposition}{Proposition}
\newcommand{\R}{\mathbb R}
\title{Two matrices with identical central zonotopal algebras\newline
but different projective configurations}
\author{Principia Math in Progress}
\date{27 September 2026}
\begin{document}
\maketitle
\begin{abstract}
The central matrix-algebra conjecture in MathDB 338791 is false.
Two explicit bridge-free $2\times4$ integer matrices have identical central
zonotopal ideals, but their column configurations have disjoint cross-ratio
orbits and hence are not $z$-equivalent.
\end{abstract}

\section*{Task 1. The original conjecture}
For a full-row-rank real matrix $A$ with columns $v_1,\ldots,v_m$, put
$Z_A=\sum_i[0,v_i]$. For each facet $H$ of $Z_A$, choose a nonzero normal
$\eta_H$, and let $m(H)$ count the nonzero entries of $\eta_H A$.
The central zonotopal algebra is
\[
 C_A^{\mathcal C}=\R[x_1,\ldots,x_n]/I_A^{\mathcal C},\qquad
 I_A^{\mathcal C}=\langle(\eta_H\cdot x)^{m(H)}:H\text{ a facet}\rangle.
\]
Two matrices are $z$-equivalent if their row spaces differ by an invertible
diagonal coordinate change and a coordinate permutation. A bridge-column
is a column whose deletion lowers the rank. The conjecture asserts that
isomorphism of central algebras, graded or ungraded, is equivalent to
$z$-equivalence after removing bridge-columns and corresponding dependent
rows. This is Conjecture 2 in the preprint and Conjecture 16 in the final
article \cite{N}. The following examples have no bridge-columns, so no
reduction is needed.

\begin{proposition}
For $t\in\R\setminus\{0,1\}$ let
\[
 A_t=\begin{pmatrix}1&0&1&1\\0&1&1&t\end{pmatrix}.
\]
Then $I_{A_t}^{\mathcal C}=(x,y)^3$. Nevertheless, $A_2$ and $A_3$ are not
$z$-equivalent. Thus they disprove the conjecture.
\end{proposition}

\begin{proof}
Every two columns of $A_t$ are independent. Hence its rank is two and
deleting any single column leaves rank two. Its zonotope has four pairs
of parallel edges, parallel to its four distinct column directions.
Indeed, for any normal $\eta$, the face maximizing $\eta\cdot z$ fixes an
endpoint of each segment with $\eta\cdot v_i\ne0$, and leaves free the
segments with $\eta\cdot v_i=0$. Such a face is an edge exactly when one
column direction is perpendicular to $\eta$.

We may therefore use the four normal directions
$(1,0),(0,1),(1,-1),(t,-1)$. Each normal annihilates exactly one column,
so all four exponents are three. Consequently
\begin{equation}\label{ideal}
 I_{A_t}^{\mathcal C}
 =\langle x^3,y^3,(x-y)^3,(tx-y)^3\rangle.
\end{equation}
Modulo $\langle x^3,y^3\rangle$, the last two generators have coefficient
matrix, in the basis $x^2y,xy^2$,
\[
 \begin{pmatrix}-3&3\\-3t^2&3t\end{pmatrix},
 \qquad\det=9t(t-1)\ne0.
\]
They therefore generate both mixed cubic monomials as well. All four
generators in \eqref{ideal} are homogeneous cubics, proving equality with
$(x,y)^3$. In particular both algebras are literally
\[
 \R[x,y]/(x,y)^3,
\]
with graded basis $1,x,y,x^2,xy,y^2$ and Hilbert series $1+2q+3q^2$.

It remains to rule out $z$-equivalence. For four nonproportional vectors
$w_1,\ldots,w_4\in\R^2$, define
\[
 r(w_1,w_2,w_3,w_4)
 =\frac{\det(w_1,w_3)\det(w_2,w_4)}
 {\det(w_1,w_4)\det(w_2,w_3)}.
\]
This number is unchanged by a common invertible linear transformation or
by independent nonzero rescalings of the vectors: the determinant factors
cancel. Permuting the vectors gives the set
\begin{equation}\label{orbit}
 \left\{r,\frac1r,1-r,\frac1{1-r},
 \frac r{r-1},\frac{r-1}r\right\}.
\end{equation}
For completeness, adjacent transpositions change the ratio by $r\mapsto
1/r$ or $r\mapsto1-r$; these generate precisely the six displayed
expressions. Coincidences can reduce the number of distinct values.
For the columns of $A_t$ in their displayed order, the ratio is $1/t$.
Thus the two permutation orbits are
\[
 \mathcal R(A_2)=\{-1,\tfrac12,2\},\qquad
 \mathcal R(A_3)=\{-2,-\tfrac12,\tfrac13,\tfrac23,\tfrac32,3\}.
\]
They are disjoint. If the row spaces of $A_2$ and a column-permuted and
column-rescaled $A_3$ were equal, their two row bases would differ by an
invertible $2\times2$ matrix. This would give projectively equivalent
column configurations, contradicting the disjoint orbits. Both algebra
isomorphism conditions hold, while the third conjectured condition fails.
\end{proof}

The same computation gives infinitely many inequivalent realizations with
this one central algebra. In fact, for $s,t>2$, \eqref{orbit} contains just
one value greater than two, namely the parameter itself; hence
$A_s\sim_z A_t$ forces $s=t$. All these matrices represent the same uniform
matroid $U_{2,4}$. The result concerns arbitrary real matrices and does not
resolve the separate conjecture for graphical or totally unimodular
configurations.

\section*{Verification and scope}
The accompanying exact verifier checks every column deletion, enumerates
the zonotope vertices and facet directions, computes the full cubic
coefficient matrices, and evaluates all 24 column orderings. As a separate
check, it rules out all 24 projective matchings by rational linear algebra,
with self-equivalences as positive controls. Computation supports the
explicit proof; no numerical approximation is needed. The source's
displayed square dimension for the reduced matrix is inconsistent with
the stated deletion operation. Here that operation deletes nothing, so
the counterexample addresses its intended meaning without using that typo.
The proof was checked by its authoring agent, without independent peer
review or a proof assistant.

\section*{Four follow-up tasks}
These are further questions, not additional campaign credits or claims of
new open problems.
\begin{description}
\item[Task 2.] Classify the central algebras of rank-two configurations
with repeated column directions, in terms of their multiplicities and
projective parameters.
\item[Task 3.] Determine for each uniform matroid $U_{n,m}$ whether all
real realizations have the same central algebra, and describe that algebra.
\item[Task 4.] Resolve the separate graphical conjecture: does the central
algebra determine the graphical matroid after bridge deletion?
\item[Task 5.] Find a natural additional structure on the central algebra
that recovers the projective column configuration, including its cross ratios.
\end{description}

\begin{thebibliography}{1}
\bibitem{N} G. Nenashev, \emph{Classification of external Zonotopal algebras},
Electronic Journal of Combinatorics 26(1) (2019), P1.32,
\href{https://doi.org/10.37236/8299}{doi:10.37236/8299};
preprint \href{https://arxiv.org/abs/1803.09966}{arXiv:1803.09966v1}.
\end{thebibliography}
\end{document}
